\documentclass[11pt, a4paper]{article}
\usepackage[a4paper, top=2.5cm, bottom=2.5cm, left=2.5cm, right=2.5cm]{geometry}
\usepackage{fontspec}

% --- UNIVERSAL PREAMBLE BLOCK ---
\usepackage[english, bidi=basic, provide=*]{babel}
\babelprovide[import, onchar=ids fonts]{english}

% Set default font to Noto Sans
\babelfont{rm}{Noto Sans}

\usepackage{amsmath}
\usepackage{enumitem}
\setlist[itemize]{label=-}
\usepackage{booktabs}
\usepackage{xcolor}
\usepackage{listings}
\usepackage{microtype}
\usepackage{titlesec}

% Define color palette for code listings and visual accents
\definecolor{codebg}{RGB}{248, 249, 250}
\definecolor{codeframe}{RGB}{220, 224, 230}
\definecolor{primary}{RGB}{31, 78, 121}
\definecolor{secondary}{RGB}{89, 89, 89}
\definecolor{keywordcolor}{RGB}{0, 102, 204}
\definecolor{stringcolor}{RGB}{163, 21, 21}
\definecolor{commentcolor}{RGB}{0, 128, 0}

% Section styling
\titleformat{\section}{\large\bfseries\color{primary}}{\thesection}{1em}{}
\titleformat{\subsection}{\normalfont\bfseries\color{primary}}{\thesubsection}{1em}{}

% Configure listing environment for Python snippets
\lstset{
  language=Python,
  backgroundcolor=\color{codebg},
  frame=single,
  rulecolor=\color{codeframe},
  basicstyle=\small\ttfamily,
  keywordstyle=\color{keywordcolor}\bfseries,
  stringstyle=\color{stringcolor},
  commentstyle=\color{commentcolor}\itshape,
  numbers=none,
  breaklines=true,
  captionpos=b,
  showstringspaces=false
}

\usepackage{hyperref}
\hypersetup{
    colorlinks=true,
    linkcolor=primary,
    urlcolor=primary,
    pdftitle={ESR Viewer Dashboard Technical Documentation}
}

\title{\Huge\bfseries\color{primary} Technical Documentation \\[0.3em] \Large ESR Viewer --- R\&D Evaluation Dashboard}
\author{\textbf{System Engineering \& Analytical Documentation}}
\date{\today}

\begin{document}

\maketitle

\begin{abstract}
This document provides complete technical specifications for the \textbf{ESR Viewer} web platform. Built using Python, Streamlit, Pandas, Matplotlib, Seaborn, and Plotly, the platform enables interactive quantitative and qualitative analysis of Evaluation Summary Reports (ESR) across temporal, geographical, and categorical dimensions.
\end{abstract}

\vspace{1.5em}

\section{Overview \& Objectives}

The primary goal of the \textbf{ESR Viewer} application is to transform raw R\&D proposal evaluation data into actionable intelligence. Evaluation Summary Reports contain scores across multiple criteria (Blocks) as well as descriptive feedback (Issues/Remarks).

Key operational capabilities include:
\begin{itemize}
    \item Multi-dimensional filtering across Area, Action/Project Type, and Evaluation Blocks.
    \item Longitudinal analysis of score progressions over time.
    \item Benchmarking of subject areas via unified global heatmaps.
    \item Evaluation volume tracking and call breakdown.
    \item Qualitative critique density analysis and correlation modeling with numerical scores.
\end{itemize}

\section{System Architecture \& Tech Stack}

The application follows a reactive, single-page dashboard design powered by Streamlit and standard Python data science libraries:

\begin{table}[htbp]
\centering
\caption{System Component Matrix}
\vspace{0.5em}
\begin{tabular}{@{}llp{8cm}@{}}
\toprule
\textbf{Layer} & \textbf{Technology} & \textbf{Role in System} \\
\midrule
UI Framework & Streamlit 1.x & State management, sidebar layout, dynamic multi-tab presentation. \\
Data Processing & Pandas & Data loading, cleaning, aggregation, dynamic pivoting, and merging. \\
Static Rendering & Matplotlib \& Seaborn & Generation of annotated heatmaps and clean standard line plots. \\
Interactive Charts & Plotly Express \& GO & Multi-axis time series, stacked bars, and grouped issue charts. \\
\bottomrule
\end{tabular}
\end{table}

\section{Input Data Specification}

The system ingests tabular data from an Excel dataset formatted as long/tidy data (\texttt{esr\_data\_long.xlsx}). 

\begin{table}[htbp]
\centering
\caption{Data Schema Requirements}
\vspace{0.5em}
\begin{tabular}{@{}lll@{}}
\toprule
\textbf{Field Name} & \textbf{Data Type} & \textbf{Description} \\
\midrule
\texttt{Year} & String / Numeric & Year of evaluation (cast to string upon import). \\
\texttt{Area} & String & R\&D domain or scientific panel category. \\
\texttt{Type} & String & Action scheme or project type (e.g., RIA, IA, CSA). \\
\texttt{Block} & String & Evaluation criterion (e.g., Excellence, Impact). \\
\texttt{Call} & String & Specific funding call identifier. \\
\texttt{Metric} & String & Distinguishes numerical scores (\texttt{Score}) from critique types. \\
\texttt{Sum} & Numeric & Total aggregated value or frequency count for the row group. \\
\texttt{Count} & Numeric & Total evaluation instances represented in the row group. \\
\bottomrule
\end{tabular}
\end{table}

\section{Core Software Architecture \& Workflow}

\subsection{Data Ingestion \& Caching}
To optimize UI responsiveness, data loading is cached using Streamlit's native memory management decorators.

\begin{lstlisting}
file_path = "esr_data_long.xlsx"

@st.cache_data
def load_data(path):
    df = pd.read_excel(path)
    df['Year'] = df['Year'].astype(str)
    return df
\end{lstlisting}

If the source dataset is missing, the platform catches \texttt{FileNotFoundError} gracefully, rendering an error alert and terminating script execution via \texttt{st.stop()}.

\subsection{Data Partitioning \& Filter Cascade}
The dataset is split into two functional subsets:
\begin{enumerate}
    \item \textbf{Scores Dataset (\texttt{df\_scores})}: Rows where \texttt{Metric == 'Score'}.
    \item \textbf{Issues Dataset (\texttt{df\_issues})}: Rows representing qualitative observations (\texttt{Metric != 'Score'}).
\end{enumerate}

Sidebar widgets provide reactive controls:
\begin{itemize}
    \item \textbf{Area}: Single selection dropdown with default "All".
    \item \textbf{Type}: Single selection dropdown with default "All".
    \item \textbf{Block}: Multi-select widget initialized to select all available blocks.
\end{itemize}

\section{Visual Analytics Modules (Tabs Breakdown)}

The user interface partitions insights into five dedicated analytical perspectives.

\subsection{Tab 1: Trend Analysis (Temporal Evolution)}
Computes the total sum and count grouped by \texttt{Year} and \texttt{Block} to render mean score evolution curves:
\begin{equation}
\text{Average Score} = \frac{\sum \text{Sum}}{\sum \text{Count}}
\end{equation}
The output is rendered as a Seaborn multi-line plot displaying score trajectories across years for selected evaluation criteria.

\subsection{Tab 2: Strengths Map (Heatmap Matrix)}
Generates a global overview matrix comparing performance across all \texttt{Area} domains versus \texttt{Block} criteria.
\begin{itemize}
    \item \textbf{Intentional Bypass}: The user's active \texttt{Area} filter is bypassed in this view to allow context-aware benchmarking against global averages.
    \item \textbf{Visualization}: Annotated Seaborn heatmap using the \texttt{RdYlGn} diverging palette.
\end{itemize}

\subsection{Tab 3: Volume \& Call Distribution}
Presents project evaluation counts grouped by evaluation year and specific funding calls. Displayed using a Plotly interactive stacked bar chart.

\subsection{Tab 4: Non-Score Issue Analyzer}
Aggregates non-score issue occurrences across evaluation criteria blocks (\texttt{Block}) and observation types (\texttt{Metric}), providing insight into critical bottlenecks or common proposal weaknesses.

\subsection{Tab 5: Score vs. Critique Frequency Dynamics}
Analyzes the potential impact of critique density on score outcomes.

\begin{enumerate}
    \item \textbf{Unified Dual Y-Axis Chart}: Displays overall average score on the primary vertical axis (bold line) alongside normalized mean appearances of various critique keywords on the secondary vertical axis.
    \item \textbf{Statistical Correlation Heatmap}: Computes Pearson correlation coefficients ($r$) between mean scores and critique metrics:
    \begin{equation}
    r = \frac{\sum (x_i - \bar{x})(y_i - \bar{y})}{\sqrt{\sum (x_i - \bar{x})^2 \sum (y_i - \bar{y})^2}}
    \end{equation}
    Rendered on a bounded scale $[-1.0, 1.0]$ using a \texttt{coolwarm} color scheme.
    \item \textbf{Consolidated Data Table}: Formatted numerical table formatted to 3 decimal places for detailed audit.
\end{enumerate}

\section{Deployment \& Local Execution}

To launch the dashboard locally:

\begin{enumerate}
    \item Install required Python dependencies:
\begin{lstlisting}[language=bash]
pip install streamlit pandas openpyxl matplotlib seaborn plotly
\end{lstlisting}
    \item Ensure \texttt{esr\_data\_long.xlsx} and \texttt{form.jpg} are present in the working directory.
    \item Launch the Streamlit server:
\begin{lstlisting}[language=bash]
streamlit run app.py
\end{lstlisting}
\end{enumerate}

\end{document}
